% Generated from the measurement described below; included by validation_report.tex.
Requirement NF-01 asks that scalar programs run no more than 5\,\% slower than before the
vector work. The \texttt{kites-cli} of the baseline commit \texttt{7ddb6f3} and of this branch
were built in the same way (Debug, Apple Clang) and run alternately in batch mode on the
single-cycle processor. \Cref{tab:perf} gives the median CPU time (user plus system) of each
program. The pipelines could not be compared, because the baseline CLI never finishes on a
pipeline model (F22).

\begin{table}[H]
\centering\small
\caption{Scalar run time, baseline and current (median CPU time).}\label{tab:perf}
\begin{tabularx}{\textwidth}{@{}YL{1.2cm}L{2.0cm}L{2.0cm}L{1.6cm}@{}}
\toprule
Program & Runs & Baseline (ms) & Current (ms) & Change \\
\midrule
cache\_matrix\_mul\_3x3.s & 21 & 206 & 197 & -4.4\,\% \\
pipeline\_test\_code\_1.s & 21 & 37 & 39 & +7.3\,\% \\
scalar loop (40\,002 instructions) & 5 & 4647 & 4350 & -6.4\,\% \\
\bottomrule
\end{tabularx}
\end{table}

Kites writes \file{registers_dump.json} after every step. The first measurement showed scalar
programs up to 65\,\% slower, because the dump had grown by the vector registers and the vector
CSRs and was formatted value by value through stream manipulators. The dump is now built in one
string with a direct hexadecimal formatter (byte-identical for the existing sections) and writes
vector registers without leading zeros. Per instruction, the current build is now faster than
the baseline. The remaining difference on the shortest program is a one-off start-up cost of
about 3\,ms for building the vector instruction tables; its wall time is unchanged.
